% GENERATED FILE. Edit canonical lessons, then run npm run generate:books.
% canonical-source-hash: fnv1a64:c4d58671f4cd84e8
% canonical-lessons: RU-C403-rul, RU-C403-tormoz, RU-C403-tormozit, RU-C403-razgonyatsya, RU-C403-signalit

\chapter{Behind the wheel}
\label{ch:ru-behind-the-wheel}
\hlchaptermodality{403}

\begin{chapteropening}
\textbf{By the end of this chapter:} \emph{I can talk about driving: the steering wheel, the brake, braking, speeding up and sounding the horn.}

Complete the last lesson of chapter 403: I can talk about driving: the steering wheel, the brake, braking, speeding up and sounding the horn.
\end{chapteropening}

\section[\texorpdfstring{\ru{руль} (rul)}{руль (rul)}]{\ru{руль} (rul) — a steering wheel; a handlebar}

\label{lesson:RU-C403-rul}



\begin{quote}
Before the new one: say the Russian for a task, then the Russian for experience.
\end{quote}

\subsection*{You'll want to know: \ru{руль}}
\textbf{\ru{руль}} (\emph{rul}) — \textquotedblleft{}a steering wheel; a handlebar\textquotedblright{}.

\begin{quote}
\ru{Папа} \ru{держал} \ru{руль} \ru{обеими} \ru{руками}.
\end{quote}

\begin{quote}
\emph{Pápa derzhál rul obéimi rukámi.}
\end{quote}

\begin{quote}
Dad held the steering wheel with both hands.
\end{quote}

Masculine; the stress moves to the ending: \emph{\ru{руля}} (\emph{rulyá}), \emph{\ru{за} \ru{рулём}} (\emph{za rulyóm}) — at the wheel, driving.

\begin{cousinweb}
\emph{\ru{руль}} was borrowed from Dutch \emph{roer} 'a rudder' (EtymDB): the word for steering a ship came to mean steering a car or a bicycle.
\end{cousinweb}

\begin{grammarlens}[title={What you've built}]
One more word for telling, explaining and describing at B1.
\end{grammarlens}

\subsection*{Your turn}
\begin{itemize}
\raggedright
  \item \emph{Say it:} \emph{rul}
  \item \emph{Say it:} \emph{rul}, once more
  \item \emph{Say it:} a sentence of your own with \emph{rul}
  \item \emph{Recall:} say the Russian for a task, then the Russian for experience, then say \emph{rul} again
\end{itemize}

\begin{tcolorbox}[breakable,colback=teal!4,colframe=teal!35!black,title={Before you move on}]
What does \ru{руль} mean? (\textquotedblleft{}A steering wheel; a handlebar\textquotedblright{}.) Say it once more.
\end{tcolorbox}
\section[\texorpdfstring{\ru{тормоз} (tórmoz)}{тормоз (tórmoz)}]{\ru{тормоз} (tórmoz) — a brake}

\label{lesson:RU-C403-tormoz}



\begin{quote}
Before the new one: say the Russian for experience, then the Russian for a steering wheel.
\end{quote}

\subsection*{You'll want to know: \ru{тормоз}}
\textbf{\ru{тормоз}} (\emph{tórmoz}) — \textquotedblleft{}a brake\textquotedblright{}.

\begin{quote}
\ru{У} \ru{этой} \ru{старой} \ru{машины} \ru{плохие} \ru{тормоза}.
\end{quote}

\begin{quote}
\emph{U étoy stároy mashíny plokhíye tormozá.}
\end{quote}

\begin{quote}
This old car has bad brakes.
\end{quote}

Plural \emph{\ru{тормоза}}. \emph{\ru{Нажать} \ru{на} \ru{тормоз}} — to put on the brake.

\begin{cousinweb}
\emph{\ru{тормоз}} goes back to Greek \emph{tórmos} 'a socket into which a peg fits' (EtymDB). How the word reached Russian is not settled.
\end{cousinweb}

\begin{grammarlens}[title={What you've built}]
One more word for telling, explaining and describing at B1.
\end{grammarlens}

\subsection*{Your turn}
\begin{itemize}
\raggedright
  \item \emph{Say it:} \emph{tórmoz}
  \item \emph{Say it:} \emph{tórmoz}, once more
  \item \emph{Say it:} a sentence of your own with \emph{tórmoz}
  \item \emph{Recall:} say the Russian for experience, then the Russian for a steering wheel, then say \emph{tórmoz} again
\end{itemize}

\begin{tcolorbox}[breakable,colback=teal!4,colframe=teal!35!black,title={Before you move on}]
What does \ru{тормоз} mean? (\textquotedblleft{}A brake\textquotedblright{}.) Say it once more.
\end{tcolorbox}
\section[\texorpdfstring{\ru{тормозить} (tormozít)}{тормозить (tormozít)}]{\ru{тормозить} (tormozít) — to brake; to slow down}

\label{lesson:RU-C403-tormozit}



\begin{quote}
Before the new one: say the Russian for a steering wheel, then the Russian for a brake.
\end{quote}

\subsection*{You'll want to know: \ru{тормозить}}
\textbf{\ru{тормозить}} (\emph{tormozít}) — \textquotedblleft{}to brake; to slow down\textquotedblright{}.

\begin{quote}
\ru{Перед} \ru{школой} \ru{водитель} \ru{всегда} \ru{тормозит}.
\end{quote}

\begin{quote}
\emph{Péred shkóloy vodítel vsegdá tormozít.}
\end{quote}

\begin{quote}
In front of the school the driver always brakes.
\end{quote}

Imperfective; the perfective is \emph{\ru{затормозить}}. Present \emph{\ru{торможу}, \ru{тормозишь}}: \ru{з} becomes \ru{ж} in the first person.

\begin{cousinweb}
\emph{\ru{тормозить}} is built on \emph{\ru{тормоз}} 'a brake' with the verb ending \emph{-\ru{ить}}. \emph{\ru{Тормоз}} goes back to Greek \emph{tórmos} 'a socket for a peg' (EtymDB).
\end{cousinweb}

\begin{grammarlens}[title={What you've built}]
One more word for telling, explaining and describing at B1.
\end{grammarlens}

\subsection*{Your turn}
\begin{itemize}
\raggedright
  \item \emph{Say it:} \emph{tormozít}
  \item \emph{Say it:} \emph{tormozít}, once more
  \item \emph{Say it:} a sentence of your own with \emph{tormozít}
  \item \emph{Recall:} say the Russian for a steering wheel, then the Russian for a brake, then say \emph{tormozít} again
\end{itemize}

\begin{tcolorbox}[breakable,colback=teal!4,colframe=teal!35!black,title={Before you move on}]
What does \ru{тормозить} mean? (\textquotedblleft{}To brake; to slow down\textquotedblright{}.) Say it once more.
\end{tcolorbox}
\section[\texorpdfstring{\ru{разгоняться} (razgonyátsya)}{разгоняться (razgonyátsya)}]{\ru{разгоняться} (razgonyátsya) — to speed up, to gather speed}

\label{lesson:RU-C403-razgonyatsya}



\begin{quote}
Before the new one: say the Russian for a brake, then the Russian for to brake.
\end{quote}

\subsection*{You'll want to know: \ru{разгоняться}}
\textbf{\ru{разгоняться}} (\emph{razgonyátsya}) — \textquotedblleft{}to speed up, to gather speed\textquotedblright{}.

\begin{quote}
\ru{На} \ru{шоссе} \ru{машина} \ru{быстро} \ru{разгоняется}.
\end{quote}

\begin{quote}
\emph{Na shossé mashína býstro razgonyáyetsya.}
\end{quote}

\begin{quote}
On the motorway the car quickly speeds up.
\end{quote}

Imperfective; perfective \emph{\ru{разогнаться}} (\emph{\ru{разгонюсь}, \ru{разгонишься}}).

\begin{cousinweb}
\emph{\ru{разгоняться}} is built from \emph{\ru{раз}-}, \emph{\ru{гонять}} 'to drive, to chase' and \emph{-\ru{ся}}. \emph{\ru{Гонять}} is inherited from Proto-Slavic \emph{ganjati}, which EtymDB traces to the PIE root \emph{g\textsuperscript{w}\textsuperscript{h}en-} 'to strike, to drive'.
\end{cousinweb}

\begin{grammarlens}[title={What you've built}]
One more word for telling, explaining and describing at B1.
\end{grammarlens}

\subsection*{Your turn}
\begin{itemize}
\raggedright
  \item \emph{Say it:} \emph{razgonyátsya}
  \item \emph{Say it:} \emph{razgonyátsya}, once more
  \item \emph{Say it:} a sentence of your own with \emph{razgonyátsya}
  \item \emph{Recall:} say the Russian for a brake, then the Russian for to brake, then say \emph{razgonyátsya} again
\end{itemize}

\begin{tcolorbox}[breakable,colback=teal!4,colframe=teal!35!black,title={Before you move on}]
What does \ru{разгоняться} mean? (\textquotedblleft{}To speed up, to gather speed\textquotedblright{}.) Say it once more.
\end{tcolorbox}
\section[\texorpdfstring{\ru{сигналить} (signálit)}{сигналить (signálit)}]{\ru{сигналить} (signálit) — to sound the horn, to honk; to signal}

\label{lesson:RU-C403-signalit}



\begin{quote}
Before the new one: say the Russian for to brake, then the Russian for to speed up.
\end{quote}

\subsection*{You'll want to know: \ru{сигналить}}
\textbf{\ru{сигналить}} (\emph{signálit}) — \textquotedblleft{}to sound the horn, to honk; to signal\textquotedblright{}.

\begin{quote}
\ru{Возле} \ru{дома} \ru{водитель} \ru{такси} \ru{долго} \ru{сигналил}.
\end{quote}

\begin{quote}
\emph{Vózle dóma vodítel taksí dólgo signálil.}
\end{quote}

\begin{quote}
Outside the house the taxi driver kept sounding his horn.
\end{quote}

Imperfective; perfective \emph{\ru{просигналить}}. Built on \emph{\ru{сигнал}} — a signal.

\begin{cousinweb}
\emph{\ru{сигналить}} is built on \emph{\ru{сигнал}} 'a signal'. \emph{\ru{Сигнал}} is usually traced to French \emph{signal}, which EtymDB derives from Latin \emph{signālis}, from \emph{signum} 'a sign'.
\end{cousinweb}

\begin{grammarlens}[title={What you've built}]
One more word for telling, explaining and describing at B1.
\end{grammarlens}

\subsection*{Your turn}
\begin{itemize}
\raggedright
  \item \emph{Say it:} \emph{signálit}
  \item \emph{Say it:} \emph{signálit}, once more
  \item \emph{Say it:} a sentence of your own with \emph{signálit}
  \item \emph{Recall:} say the Russian for to brake, then the Russian for to speed up, then say \emph{signálit} again
\end{itemize}

\begin{tcolorbox}[breakable,colback=teal!4,colframe=teal!35!black,title={Before you move on}]
What does \ru{сигналить} mean? (\textquotedblleft{}To sound the horn, to honk; to signal\textquotedblright{}.) Say it once more.
\end{tcolorbox}
